\documentclass[11pt]{article}
\usepackage[utf8]{inputenc}
\usepackage[T1]{fontenc}
\usepackage[a4paper, margin=1in]{geometry}
\usepackage{setspace}
\usepackage{microtype}
\usepackage{lmodern}

\pagestyle{empty}
\setlength{\parindent}{0pt}

\begin{document}

\begin{center}
{\Large\bfseries The Political Economy of Stagnation: Salariat Formation,\\[2pt]
Class Location, and the Inhibition of Capital\\[2pt]
Accumulation in Assam\par}

\vspace{14pt}

Kabeer Bora\\[2pt]
\textit{University of Missouri--Kansas City}

\vspace{18pt}

\textbf{ABSTRACT}
\end{center}

\vspace{4pt}

This paper offers a long-run political-economy analysis of Assam's structural underdevelopment through the lens of heterodox class analysis. We argue that stagnation is the endogenous product of a historically entrenched salariat---a class fraction whose material reproduction depends on state patronage---formed through the articulation of precolonial corv\'ee labor relations with the demands of British colonial administration.

\vspace{6pt}

The paper makes two theoretical contributions. First, we apply and modify Pranab Bardhan's (1984) dominant-proprietary-class framework: in Assam, the industrial-capitalist and rich-farmer fractions were historically absent, producing a structurally distinct configuration of uncontested salariat hegemony. Second, we use Nicos Poulantzas (1975) and Erik Olin Wright (1985) to specify the political and ideological forms of this hegemony. Poulantzas grounds our reading of the salariat as a class fraction constituted at the intersection of the mental--manual labor division and state administrative function; Wright supplies the concept of credential-asset exploitation, by which the salariat appropriates a continuous share of state-claimed surplus through its monopoly over scarce credentials. The two together explain why threats to that monopoly---most visibly, educated immigrant competition---register as existential class dispossession and are met with the ethnic, sons-of-soil mobilizations that have defined Assam's modern political history. The political and ideological form of class struggle, not merely the economic, is essential to the analysis.

\vspace{12pt}

\textit{JEL Classification: B51, O18, N35, J45, R11}

\end{document}
